% Generated from SKILL.md by tools/build_skill_appendix.py; do not edit by hand.
English edition of the \passthrough{\lstinline!experiment-exploration!}
skill; keep the two in sync.

This approach comes from the SRCNN FPGA optimization in
\passthrough{\lstinline!\~/School/Hardware/remote/90096\_CNN\_FPGA!}:
170+ experiments, each traceable to its experiment card, raw reports,
and commit. Scripts and templates live in this skill's directory
(\passthrough{\lstinline!scripts/!},
\passthrough{\lstinline!templates/!}).

Division of labor with other skills:

\begin{itemize}
\tightlist
\item
  This skill covers the phase where you don't yet know which approach
  works. Experiment code lives on experiment branches and in independent
  copies, never on the mainline; once an approach is chosen, implement
  it properly following
  \passthrough{\lstinline!production-engineering!}.
\item
  For experiments on n08 / n09, follow
  \passthrough{\lstinline!lab-gpu-servers!} for connecting, picking
  GPUs, and moving data.
\end{itemize}

\subsection{Before starting: write the exploration
plan}

Create an exploration plan under \passthrough{\lstinline!docs/!} from
the template
\passthrough{\lstinline!templates/plan.md!}.
Fill in the following before the first experiment:

\begin{enumerate}
\def\labelenumi{\arabic{enumi}.}
\tightlist
\item
  \textbf{Problem and invariant contract}: what must not change
  (interfaces, data formats, numeric precision, test set). If these
  change, it is no longer the same problem.
\item
  \textbf{A single primary metric}: one number compares every approach,
  with the conversion formula fixed in writing
  (e.g.~\passthrough{\lstinline!ms/image = cycles × clock period (ns) / 10\^6!}).
  Numbers measured on different bases are not ranked together; when they
  must appear side by side, label the basis.
\item
  \textbf{Gates}, in three kinds:

  \begin{itemize}
  \tightlist
  \item
    Quality gate: how correct the result must be (e.g.~\(\Delta\)PSNR
    \(\ge\) \$-\$0.1 dB, bit-exact against the reference);
  \item
    Screening gate: whether it is worth taking to the next level of
    verification;
  \item
    Final gate: whether it can be declared ``feasible''. Only approaches
    that pass the final gate enter the ranking.
  \end{itemize}
\item
  \textbf{Evidence levels}: list your domain's evidence from weakest to
  strongest, e.g.~``paper estimate \textless{} local / unit experiment
  \textless{} simulation \textless{} synthesis estimate \textless{}
  implementation \textless{} hardware measurement''. Tag every
  conclusion with the level it has reached; never present lower-level
  evidence as higher-level.
\item
  \textbf{Numbering scheme}: series prefixes by direction (e.g.~E for
  baseline and sweeps, P / W / G for individual optimization directions,
  C for combinations, F for the mainline); IDs are never reused; reserve
  a separate range for replays (e.g.~100 and above).
\item
  \textbf{Budget and stop rules}: timeout per run; maximum time on any
  one problem (e.g.~45 minutes), after which you stop and record;
  concurrency and memory rules for the machines.
\item
  \textbf{Protected artifacts}: baseline source, data, reference
  outputs. Record their sha256 before starting and verify again at
  wrap-up.
\item
  \textbf{User decisions}: record them in the plan with dates. When a
  new decision supersedes an old rule, keep the old rule as history;
  don't rewrite it.
\end{enumerate}

\subsection{Lifecycle of one
experiment}

\begin{longtable}[]{@{}
  >{\raggedright\arraybackslash}p{(\linewidth - 2\tabcolsep) * \real{0.28}}
  >{\raggedright\arraybackslash}p{(\linewidth - 2\tabcolsep) * \real{0.72}}@{}}
\toprule\noalign{}
\begin{minipage}[b]{\linewidth}\raggedright
Step
\end{minipage} & \begin{minipage}[b]{\linewidth}\raggedright
What to do
\end{minipage} \\
\midrule\noalign{}
\endhead
\bottomrule\noalign{}
\endlastfoot
1. Assign an ID & Confirm it is unused. Completed IDs are never reused;
a rerun after a fix also gets a new ID, linked back with a
\passthrough{\lstinline!retry!} relation \\
2. Write and commit the experiment card &
\passthrough{\lstinline!templates/experiment.json!}:
purpose, control experiment, \textbf{the only thing changed} relative to
the control, conditions held fixed, paper estimate and its basis, how it
will be judged, and what each possible outcome leads to next.
\textbf{Commit first, then run} \\
3. Run in isolation & Copy an independent design or codebase from the
control; don't touch files of completed experiments; no concurrent runs
in the same working directory; every run has a timeout and a log \\
4. Archive & \passthrough{\lstinline!experiments/<id>\_<short-name>/!}:
a copy of the card, raw logs, raw reports, parsed metrics, a one-page
note (result, deviation from the estimate, conclusion), plus the source
commit and tool versions \\
5. Register & The registry \passthrough{\lstinline!results.csv!} is
appended by the parsing script, never by hand; leave unmeasured values
blank, not 0; failures, timeouts, and invalid configurations get a row
too \\
6. Record & Append an entry to \passthrough{\lstinline!LOG.md!}; update
the current state in \passthrough{\lstinline!PROGRESS.md!} \\
7. Commit & One commit per experiment, message starting with the ID,
e.g.~\passthrough{\lstinline!F25 nested\_weight\_loader!} \\
8. Update the evolution graph & Add the node and relations to
\passthrough{\lstinline!graph.json!} and regenerate the SVGs (section
6) \\
\end{longtable}

Completed experiments are immutable: to reproduce one, assign a new ID,
copy the archived design, and rerun it with the recorded commands.

\subsection{When proposing an
optimization}

\begin{enumerate}
\def\labelenumi{\arabic{enumi}.}
\tightlist
\item
  \textbf{Find the bottleneck first, with evidence}: a specific line
  from a report or profile, e.g.~``the 2.822 ns global critical path
  sits on the 15-bit address multiply in Conv1 weight loading''. No
  experiment is opened on an evidence-free ``this should be faster''.
\item
  \textbf{State a testable hypothesis}: ``Changing X should move the
  metric from a to b, because \ldots; the likely cost is \ldots{}''.
\item
  \textbf{Pick a control and change only this one thing}: name the
  control experiment and write the difference in the card's
  \passthrough{\lstinline!changes!}. Change several factors at once and
  you can't tell which one mattered.
\item
  \textbf{Commit the paper estimate first}: write the estimate and its
  basis into the card and commit before running, so it can't be revised
  after the fact.
\item
  \textbf{Compare the measurement with the estimate}: when they differ a
  lot, first explain why (where the estimate went wrong), then decide
  the next step.
\item
  \textbf{Turn the conclusion into a next step}: adopt as the new
  baseline, keep as a branch, or stop, and say why.
\item
  \textbf{Update the evolution graph}.
\end{enumerate}

Rules of judgment:

\begin{itemize}
\tightlist
\item
  Measure the baseline and fix the numeric contract before optimizing.
\item
  Sweep one parameter dimension at a time, round by round; don't change
  every parameter at once.
\item
  Local results (a unit, a single layer, a small kernel) don't
  extrapolate to the whole system.
\item
  Gains from two approaches don't simply add up; a combination must be
  measured on its own.
\item
  The fastest point that fails a gate stays out of the ranking, e.g.~a
  combination with the lowest latency but three times the LUT budget.
\item
  Failure is a result: keep the node and state the cause. Distinguish
  ``configuration didn't take effect'' (invalid: fix and retry) from
  ``valid but below the bar'' (a conclusion).
\item
  A tool ``exiting successfully'' doesn't mean the result is correct;
  check the actual metrics.
\item
  When one problem hits its budget, stop, record the current state, and
  switch direction or ask the user.
\end{itemize}

\subsection{File layout}

\begin{lstlisting}
docs/NNNN-<topic>.md            exploration plan; conclusions written at wrap-up
exp/configs/<id>.json           experiment cards (committed before running)
experiments/
  PROGRESS.md                   resume entry point: current state, what's running, next steps, how to resume
  LOG.md                        append-only event log
  results.csv                   registry (script-generated)
  graph.json                    evolution graph data: hand-written ideas and relations
  evolution/                    generated .svg / .mmd (don't edit by hand)
  <id>_<short-name>/            per-experiment archive
\end{lstlisting}

\passthrough{\lstinline!PROGRESS.md!} template:
\passthrough{\lstinline!templates/PROGRESS.md!}.

\subsection{Registry and log}

\textbf{\passthrough{\lstinline!results.csv!}}: one row per experiment,
with at least these columns:

\passthrough{\lstinline!id, time, series, control, changes, status, <primary metric>, <gate-related metrics>, evidence\_level, git\_commit, archive, notes!}

\begin{itemize}
\tightlist
\item
  \passthrough{\lstinline!status!} uses the same values as the evolution
  graph:
  \passthrough{\lstinline!pass / warn / fail / invalid / timeout / planned / analysis / partial!}.
\item
  Predictions stay in the experiment card, never in the registry; the
  registry holds measurements only.
\item
  If a registered result turns out to be wrong, append a correction
  event to LOG.md and fix it in a new experiment; don't rewrite the
  original reports and logs.
\end{itemize}

\textbf{\passthrough{\lstinline!LOG.md!}} is append-only, one line per
entry, with UTC time and IDs:

\begin{lstlisting}
- 2026-10-10T16:52:16Z: E023/E024/E025 invalid: syn.directive silently ignored by the tool. Now generating directives from the schema and rejecting unknown options; retried as E028/E029/E032.
\end{lstlisting}

\subsection{Experiment evolution graph (SVG directed
graph)}

The evolution graph answers three questions: \textbf{what each
experiment builds on, why it was adjusted that way, and how its result
led to the next experiment}. Update it when a new optimization direction
opens, when a phase wraps up, and when reporting to the user.

\textbf{Data}: \passthrough{\lstinline!experiments/graph.json!}.
Hand-write a one-line idea for each experiment and where it came from.
Example:
\passthrough{\lstinline!templates/graph.example.json!};
rendered result:
\passthrough{\lstinline!templates/example-evolution/overview.svg!}.

\begin{lstlisting}
{ "id": "F25", "title": "Disable load-loop flattening", "group": "mainline", "status": "warn",
  "idea": "One change only: LOOP_FLATTEN off", "metric": 2.64, "problem": "Timing warning at 2.5 ns",
  "parents": ["F24"] }
\end{lstlisting}

\textbf{Generate} (requires Node and dagre 0.8.x; see the header of the
script):

\begin{lstlisting}[language=bash]
node <this skill>/scripts/evolution_graph.cjs experiments/graph.json --all \
     --results experiments/results.csv --metric-column latency_ms
node <this skill>/scripts/evolution_graph.cjs experiments/graph.json --all \
     --results experiments/results.csv --metric-column latency_ms --check
\end{lstlisting}

\passthrough{\lstinline!--all!} produces the overview
\passthrough{\lstinline!overview.svg!} plus one graph per route
(\passthrough{\lstinline!group!}); route graphs draw direct sources from
outside the route as grey dashed boxes.
\passthrough{\lstinline!--check!} writes nothing; it only checks the
data rules, node overlap, and whether the generated files are stale.

\textbf{Node statuses}:

\begin{longtable}[]{@{}
  >{\raggedright\arraybackslash}p{(\linewidth - 4\tabcolsep) * \real{0.3333}}
  >{\raggedright\arraybackslash}p{(\linewidth - 4\tabcolsep) * \real{0.3333}}
  >{\raggedright\arraybackslash}p{(\linewidth - 4\tabcolsep) * \real{0.3333}}@{}}
\toprule\noalign{}
\begin{minipage}[b]{\linewidth}\raggedright
status
\end{minipage} & \begin{minipage}[b]{\linewidth}\raggedright
Meaning
\end{minipage} & \begin{minipage}[b]{\linewidth}\raggedright
Requirement
\end{minipage} \\
\midrule\noalign{}
\endhead
\bottomrule\noalign{}
\endlastfoot
\passthrough{\lstinline!pass!} & Passed the current stage, with a
comparable metric & Must have a metric; may add
\passthrough{\lstinline!"selected": true!} to mark the current choice
(\(\star\), bold border) \\
\passthrough{\lstinline!warn!} & Passed with warnings (e.g.~timing
warnings) & Must have a metric; put the warning in
\passthrough{\lstinline!problem!} \\
\passthrough{\lstinline!fail!} & Below the bar, or failed & Must have a
one-line \passthrough{\lstinline!problem!}; any metric is shown only in
grey \\
\passthrough{\lstinline!invalid!} & Invalid configuration (didn't take
effect); faded & Must have \passthrough{\lstinline!problem!} \\
\passthrough{\lstinline!timeout!} & Timed out or interrupted; faded &
Must have \passthrough{\lstinline!problem!} \\
\passthrough{\lstinline!planned!} & Planned, not run; dashed border &
Must not have a metric \\
\passthrough{\lstinline!analysis!} & Software analysis or numeric study
& --- \\
\passthrough{\lstinline!partial!} & Local experiment (unit, single
layer), no whole-system metric & Don't invent a whole-system metric \\
\end{longtable}

\textbf{Relations} (solid lines are design evolution; dashed lines are
evidence and verification):

\begin{longtable}[]{@{}lll@{}}
\toprule\noalign{}
kind & Meaning & Line \\
\midrule\noalign{}
\endhead
\bottomrule\noalign{}
\endlastfoot
\passthrough{\lstinline!source!} & Design or parameter baseline
(default) & solid \\
\passthrough{\lstinline!retry!} & Retry after a fix & solid \\
\passthrough{\lstinline!integrate!} & Integrates an existing module &
solid \\
\passthrough{\lstinline!reference!} & Numeric basis or inspiration &
dashed \\
\passthrough{\lstinline!control!} & Control & dashed \\
\passthrough{\lstinline!verify!} & Verification of the same design &
dashed \\
\passthrough{\lstinline!replay!} & Reproduction & dashed \\
\end{longtable}

\textbf{Drawing rules}:

\begin{itemize}
\tightlist
\item
  Exactly one node per experiment; relations are dependencies between
  ideas, not Git commit history; no cycles.
\item
  An experiment normally links to one primary source; link several only
  when integrating, or when bringing in a key numeric basis.
\item
  Keep failed nodes, with a specific one-line problem such as ``LUT
  248,797, 3\(\times\) over budget'', not just ``failed''.
\item
  Put only the primary metric on a node; resources, accuracy, and other
  details go in the experiment record.
\item
  The registry is the source of truth for metrics: cross-check with
  \passthrough{\lstinline!--results!}, and fix
  \passthrough{\lstinline!graph.json!} when they disagree.
\item
  Once there are many experiments (dozens or more), split them into
  routes by \passthrough{\lstinline!group!} and keep only the trunk in
  the overview; if needed, use route-level summary nodes there.
\end{itemize}

\textbf{Accompanying text}: in the exploration plan or a separate
document, write one paragraph per route answering the three questions
above, and embed the matching SVG. Readers look at the overview first,
then at each route for detail.

\subsection{Resuming and handover}

\begin{itemize}
\tightlist
\item
  \passthrough{\lstinline!PROGRESS.md!} is the resume entry point:
  current branch and baseline, completed IDs, running jobs (machine,
  process, log path), to-dos, user decisions in force, and resume steps.
  Update it at the end of every experiment.
\item
  Check processes before resuming: remote jobs may have finished or may
  still be running. For finished jobs, only collect results
  (collect-only) and don't relaunch them; never start a second copy of a
  job that is still running.
\item
  The handover document states which conclusions have reached which
  evidence level, what hasn't been verified yet, and the candidates for
  the next step with reasons.
\end{itemize}

\subsection{Reporting}

Every report uses a table; write \passthrough{\lstinline!unknown!} for
anything not measured:

\begin{longtable}[]{@{}
  >{\raggedright\arraybackslash}p{(\linewidth - 12\tabcolsep) * \real{0.1429}}
  >{\raggedright\arraybackslash}p{(\linewidth - 12\tabcolsep) * \real{0.1429}}
  >{\raggedright\arraybackslash}p{(\linewidth - 12\tabcolsep) * \real{0.1429}}
  >{\raggedright\arraybackslash}p{(\linewidth - 12\tabcolsep) * \real{0.1429}}
  >{\raggedright\arraybackslash}p{(\linewidth - 12\tabcolsep) * \real{0.1429}}
  >{\raggedright\arraybackslash}p{(\linewidth - 12\tabcolsep) * \real{0.1429}}
  >{\raggedright\arraybackslash}p{(\linewidth - 12\tabcolsep) * \real{0.1429}}@{}}
\toprule\noalign{}
\begin{minipage}[b]{\linewidth}\raggedright
ID
\end{minipage} & \begin{minipage}[b]{\linewidth}\raggedright
What changed vs.~control
\end{minipage} & \begin{minipage}[b]{\linewidth}\raggedright
Primary metric
\end{minipage} & \begin{minipage}[b]{\linewidth}\raggedright
Gate results
\end{minipage} & \begin{minipage}[b]{\linewidth}\raggedright
Evidence level
\end{minipage} & \begin{minipage}[b]{\linewidth}\raggedright
Status
\end{minipage} & \begin{minipage}[b]{\linewidth}\raggedright
Conclusion / next step
\end{minipage} \\
\midrule\noalign{}
\endhead
\bottomrule\noalign{}
\endlastfoot
\end{longtable}

Also state:

\begin{itemize}
\tightlist
\item
  the currently selected approach, and why it is the one;
\item
  deviations from the paper estimates;
\item
  items that haven't reached the final gate yet;
\item
  the path to the evolution graph.
\end{itemize}

\subsection{Subagents and models}

Subagents default to Sonnet; for the setting and the precedence order,
see section 7 of \passthrough{\lstinline!production-engineering!}.

\begin{itemize}
\tightlist
\item
  Good fits for Sonnet subagents: running an experiment from its
  committed experiment card, checking remote jobs and collecting
  results, parsing reports and appending to the registry, generating and
  checking the evolution graph, searching past experiments and logs.
\item
  Pass \passthrough{\lstinline!model: "opus"!} explicitly for: finding
  bottleneck evidence in reports, writing hypotheses and paper
  estimates, explaining gaps between estimates and measurements,
  verifying any claim that something passes the final gate.
\end{itemize}
